\section{Los dos grandes cuellos de botella de los agentes 2B}

\subsection{Entorno y despliegue}

Yao Shunyu (姚顺雨) señala que, aunque hoy se detuviera por completo el entrenamiento de modelos, con solo desplegar los modelos existentes en empresas de todo el mundo podría lograrse un impacto de 5\%--10\% en el PIB. Pero hoy el impacto real de la IA en el PIB está muy por debajo de 1\%. Esto muestra que el cuello de botella no está en el modelo, sino en el \textbf{despliegue y el entorno}.

\subsection{La urgencia de la educación}

\begin{importantbox}{La brecha entre las personas se está ampliando}
No es que la IA esté sustituyendo el trabajo de las personas, sino que \textbf{las personas que saben usar herramientas de IA están sustituyendo a las que no saben usarlas}. Es como cuando se inventó la computadora: entre quienes aprendieron a programar frente a quienes seguían usando el ábaco, la diferencia era enorme. Una de las cosas más significativas que China puede hacer hoy es una mejor educación en IA.
\end{importantbox}

